% =============================================================================
% Appendix B: User Manual
% =============================================================================
\chapter{User Manual}
\label{app:user-manual}

This appendix provides step-by-step instructions for installing, configuring, and running the SciFig-Evaluation system. All commands assume a Unix-like environment (macOS or Linux).

% -----------------------------------------------------------------------------
\section{System Requirements}
\label{sec:system-requirements}

\begin{itemize}
  \item \textbf{Python} 3.11 or later (tested on 3.12)
  \item \textbf{Node.js} 20+ and \texttt{npm} 10+ (for the dashboard)
  \item \textbf{API credentials} (at least one of):
    \begin{itemize}
      \item OpenRouter API key (\texttt{OPENROUTER\_API\_KEY})
      \item Azure OpenAI endpoint and key (\texttt{AZURE\_OPENAI\_ENDPOINT}, \texttt{AZURE\_OPENAI\_API\_KEY})
      \item Google Cloud project with Vertex AI enabled (\texttt{GCP\_PROJECT})
    \end{itemize}
  \item \textbf{Key Python packages}: \texttt{openai}, \texttt{google-genai}, \texttt{python-dotenv}, \texttt{numpy}, \texttt{matplotlib}, \texttt{scipy}
  \item Approximately 2\,GB disk space for the full dataset and outputs
\end{itemize}

% -----------------------------------------------------------------------------
\section{Installation and Setup}
\label{sec:installation}

\begin{enumerate}
  \item \textbf{Clone the repository}:
\begin{verbatim}
  git clone <repo-url> SciFig-Evaluation
  cd SciFig-Evaluation
\end{verbatim}

  \item \textbf{Create a virtual environment and install dependencies}:
\begin{verbatim}
  python3 -m venv .venv
  source .venv/bin/activate
  pip install openai google-genai python-dotenv numpy
  pip install matplotlib scipy seaborn
\end{verbatim}

  \item \textbf{Configure environment variables}. Create a \texttt{.env} file in the project root:
{\scriptsize
\begin{verbatim}
  # OpenRouter (for Llama, Gemma, Qwen, Claude, Gemini)
  OPENROUTER_API_KEY=sk-or-...

  # Azure OpenAI (for GPT-5.2, Phi-4-Multimodal)
  AZURE_OPENAI_ENDPOINT=https://your-resource.openai.azure.com/
  AZURE_OPENAI_API_KEY=...

  # Google Cloud / Vertex AI (for Gemini 3.1 Pro)
  GCP_PROJECT=your-gcp-project-id
  GOOGLE_CLOUD_LOCATION=us-central1
\end{verbatim}
}

  \item \textbf{Verify the dataset} is in place:
\begin{verbatim}
  ls Dataset/figures/english_only/   # PNG figure images
  ls Dataset/groundtruth/english_only/ # JSON annotations
\end{verbatim}
\end{enumerate}

% -----------------------------------------------------------------------------
\section{Running the Evaluation Pipeline}
\label{sec:running-pipeline}

The evaluation pipeline has three stages: generation, MQM evaluation, and results computation. Figure~\ref{fig:pipeline-flow} illustrates the overall data flow.

\begin{figure}[ht]
\centering
\begin{tikzpicture}[
  node distance=1.4cm,
  box/.style={rectangle, draw, rounded corners, minimum width=3.8cm, minimum height=0.8cm, align=center, font=\small},
  arrow/.style={-{Stealth[length=3mm]}, thick}
]
  \node[box, fill=blue!10] (dataset) {Dataset/\\figures + groundtruth};
  \node[box, fill=green!10, right=2.0cm of dataset] (gen) {Stage 1:\\LLM Generation};
  \node[box, fill=orange!10, right=2.0cm of gen] (eval) {Stage 2:\\MQM Evaluation};
  \node[box, fill=red!10, below=1.2cm of eval] (results) {Stage 3:\\Compute Results};
  \node[box, fill=purple!10, below=1.2cm of dataset] (dash) {Dashboard\\Visualisation};

  \draw[arrow] (dataset) -- (gen);
  \draw[arrow] (gen) -- (eval);
  \draw[arrow] (eval) -- (results);
  \draw[arrow] (results) -| (dash);
  \draw[arrow, dashed] (dataset) -- (dash);
\end{tikzpicture}
\caption{Evaluation pipeline data flow.}
\label{fig:pipeline-flow}
\end{figure}

\subsection{Stage 1: Generate Model Descriptions}

Each model processes all 1,005 figures across five language subfolders. The runner script automatically skips figures that have already been generated.

{\scriptsize
\begin{verbatim}
  # Generate descriptions for a single model
  python3 scripts/llm_generation/run.py gpt-5.2 --workers 4

  # Generate for a specific language subset
  python3 scripts/llm_generation/run.py gemini-3.1-pro --subfolder english_only

  # Output is written to: output/generation/<model>/<subfolder>/<fig_key>.json
\end{verbatim}
}

Available models include: \texttt{gpt-5.2}, \texttt{phi-4-multimodal}, \texttt{gemini-3.1-pro}, \texttt{llama4-maverick}, \texttt{llama4-scout}, \texttt{claude-opus-4.6}, \texttt{qwen3-vl-8b}, \texttt{qwen3-vl-30b-a3b}, \texttt{qwen3-vl-32b}, \texttt{qwen3-vl-235b-a22b}, \texttt{gemma3-4b-it}, \texttt{gemma3-12b-it}, \texttt{gemma3-27b-it}.

\subsection{Stage 2: Run MQM Evaluation}

Three judge configurations are available: reference-free (Judge~A), reference-only (Judge~B), and reference-with-image (Judge~C).

{\scriptsize
\begin{verbatim}
  # Reference-free evaluation (image + prompt, no groundtruth)
  python3 scripts/evaluation/evaluate_reference_free.py gpt-5.2 \
      --judge azure/gpt-4.1 --workers 4

  # Reference-only evaluation (groundtruth, no image)
  python3 scripts/evaluation/evaluate_reference_only.py gpt-5.2 \
      --judge azure/gpt-4.1 --workers 4

  # Reference-with-image evaluation (image + groundtruth)
  python3 scripts/evaluation/evaluate_reference_with_image.py gpt-5.2 \
      --judge azure/gpt-4.1 --workers 4

  # Output: output/evaluation/<judge_type>/<judge_slug>/<model>/<subfolder>/<fig>.json
\end{verbatim}
}

\subsection{Stage 3: Compute Results}

Scripts in \texttt{scripts/evaluation-results/} compute the tables and plots for each research question:

{\scriptsize
\begin{verbatim}
  # RQ1: Leaderboard with bootstrap CIs
  python3 scripts/evaluation-results/rq1/compute_leaderboard.py

  # RQ1: Error analysis breakdown
  python3 scripts/evaluation-results/rq1/compute_error_analysis.py

  # RQ2: Degradation under transforms
  python3 scripts/evaluation-results/rq2/compute_degradation.py

  # RQ3: Capability question accuracy
  python3 scripts/evaluation-results/rq3/compute_capability.py

  # RQ4: Agreement reliability (Krippendorff's alpha)
  python3 scripts/evaluation-results/rq4/compute_ari.py
\end{verbatim}
}

% -----------------------------------------------------------------------------
\section{Running Adversarial Experiments}
\label{sec:adversarial-experiments}

The adversarial evaluation suite operates on a curated 45-figure subset stored in \texttt{adversarial\_experiments/}. Each experiment type has a dedicated runner and evaluator script.

\begin{description}
  \item[Image Transforms] Apply visual degradations (JPEG compression, noise, rotation, blur, etc.) and regenerate descriptions:
{\scriptsize
\begin{verbatim}
  python3 scripts/experiments/run_transforms.py gpt-5.2 --workers 4
  python3 scripts/experiments/run_transforms.py gpt-5.2 --transform selective_blur
\end{verbatim}
}

  \item[Caption Bias] Test whether models resist poisoned captions or depend on caption text:
{\scriptsize
\begin{verbatim}
  python3 scripts/experiments/run_caption_bias.py gpt-5.2 --condition modified_caption
  python3 scripts/experiments/run_caption_bias.py gpt-5.2 --condition image_only
\end{verbatim}
}

  \item[Prompt Reverse] Send true/false statements about a figure; the model must classify:
{\scriptsize
\begin{verbatim}
  python3 scripts/experiments/run_prompt_reverse.py gpt-5.2 --workers 4
\end{verbatim}
}

  \item[Capability Questions] Send 5 questions per prompt (capability, hallucination, prompt reverse):
{\scriptsize
\begin{verbatim}
  python3 scripts/experiments/run_questions.py gpt-5.2 --prompt 1 --workers 4
  python3 scripts/experiments/run_questions.py gpt-5.2 --prompt 2 --workers 4
\end{verbatim}
}

  \item[Inductance Probes] Deep reasoning questions on original clean images:
{\scriptsize
\begin{verbatim}
  python3 scripts/experiments/run_inductance.py --model gemini-3.1-pro
\end{verbatim}
}

  \item[Active Admittance] Reasoning questions on selectively blurred images where the answer resolves to the blurred element:
{\scriptsize
\begin{verbatim}
  python3 scripts/experiments/run_active_admittance.py --model gpt-5.2 --workers 4
\end{verbatim}
}
\end{description}

After generation, evaluate adversarial outputs with:
{\scriptsize
\begin{verbatim}
  python3 scripts/evaluation/evaluate_transforms.py gpt-5.2 --judge azure/gpt-4.1
  python3 scripts/experiments/evaluate_caption_bias.py --model gpt-5.2
  python3 scripts/experiments/evaluate_prompt_reverse.py --model gpt-5.2
  python3 scripts/experiments/evaluate_capability.py --model gpt-5.2
  python3 scripts/experiments/evaluate_admittance.py --model gpt-5.2
\end{verbatim}
}

% -----------------------------------------------------------------------------
\section{Launching the Dashboard}
\label{sec:dashboard}

The interactive dashboard is a React application built with Vite, Tailwind CSS~4, and Recharts.

{\scriptsize
\begin{verbatim}
  cd dashboard/
  npm install
  npm run dev          # starts on http://localhost:5173
\end{verbatim}
}

Before launching, export evaluation data into the dashboard's public folder:
{\scriptsize
\begin{verbatim}
  python3 scripts/evaluation/export_dashboard_data.py --judge gpt-4o --judge gpt-4o-mini
\end{verbatim}
}

The dashboard provides the following pages:
\begin{itemize}
  \item \textbf{Dataset Browser} -- browse all figures with metadata and groundtruth
  \item \textbf{Judge Review} -- inspect MQM error annotations with highlighted text spans
  \item \textbf{Analytics} -- aggregate statistics, score distributions, error breakdowns
  \item \textbf{Final Results} -- leaderboard tables and cross-language heatmaps
  \item \textbf{Adversarial Results} -- transform degradation charts and experiment outcomes
  \item \textbf{Transform Evaluation} -- side-by-side comparison of original vs.\ transformed descriptions
  \item \textbf{Adversarial Sample Browser} -- detailed per-figure adversarial probe results
\end{itemize}

% -----------------------------------------------------------------------------
\section{Reproducing Thesis Results}
\label{sec:reproducing-results}

Table~\ref{tab:result-scripts} maps each thesis artefact to the script that produces it.

\begin{table}[ht]
\centering
\caption{Mapping of thesis tables and figures to computation scripts.}
\label{tab:result-scripts}
{\small
\begin{tabular}{@{}lll@{}}
\toprule
\textbf{Artefact} & \textbf{Script path} \\
\midrule
Leaderboard table (RQ1) & \texttt{evaluation-results/rq1/compute\_leaderboard.py} \\
Error analysis (RQ1) & \texttt{evaluation-results/rq1/compute\_error\_analysis.py} \\
Cross-lingual comparison & \texttt{evaluation-results/rq1/compute\_crosslingual.py} \\
Heatmap plot & \texttt{evaluation-results/rq1/plot\_heatmap.py} \\
Dumbbell plot & \texttt{evaluation-results/rq1/plot\_leaderboard\_dumbbell.py} \\
Scaling plot & \texttt{evaluation-results/rq1/plot\_scaling.py} \\
Human--judge correlation & \texttt{evaluation-results/rq1/compute\_human\_judge\_correlation.py} \\
Scatter plot (human vs.\ judge) & \texttt{evaluation-results/rq1/plot\_scatter\_human\_judge.py} \\
Degradation slopes (RQ2) & \texttt{evaluation-results/rq2/compute\_degradation.py} \\
Degradation slope plot & \texttt{evaluation-results/rq2/plot\_degradation\_slope.py} \\
Capability accuracy (RQ3) & \texttt{evaluation-results/rq3/compute\_capability.py} \\
Capability grouped bars & \texttt{evaluation-results/rq3/plot\_capability\_grouped\_bar.py} \\
ARI scores (RQ4) & \texttt{evaluation-results/rq4/compute\_ari.py} \\
ARI radar/dot-strip plots & \texttt{evaluation-results/rq4/plot\_ari\_radar.py} \\
Admittance stacked bars & \texttt{evaluation-results/rq4/plot\_admittance\_stacked.py} \\
Rank flow (Discussion) & \texttt{evaluation-results/discussion/plot\_rank\_flow.py} \\
Transform gallery & \texttt{evaluation-results/plot\_transforms\_gallery.py} \\
\bottomrule
\end{tabular}
}
\end{table}

All script paths above are relative to \texttt{scripts/}. Run each with:
{\scriptsize
\begin{verbatim}
  python3 scripts/evaluation-results/rq1/compute_leaderboard.py
\end{verbatim}
}

Output artefacts (JSON tables and PDF/PNG plots) are written to \texttt{output/evaluation-results/}.


% =============================================================================
% Appendix C: Maintenance Manual
% =============================================================================
\chapter{Maintenance Manual}
\label{app:maintenance-manual}

This appendix documents the internal architecture of SciFig-Evaluation to support ongoing development and extension.

% -----------------------------------------------------------------------------
\section{Repository Structure}
\label{sec:repo-structure}

{\scriptsize
\begin{verbatim}
SciFig-Evaluation/
|-- Dataset/
|   |-- figures/                 # PNG images (5 subfolders by language)
|   |   |-- english_only/
|   |   |-- german_only/
|   |   |-- bulgarian_only/
|   |   |-- chinese_only/
|   |   +-- multi_language/
|   +-- groundtruth/             # Human-authored JSON annotations (same subfolders)
|-- adversarial_experiments/
|   |-- samples.json             # 45-figure sample manifest
|   |-- figures/                 # Original + transformed images per figure
|   |-- benchmarks/              # Question/probe JSONs for each experiment
|   +-- papers/                  # Source PDFs for adversarial figures
|-- atomic_mqm/
|   |-- atoms/                   # 120 atomic MQM decomposition JSONs
|   |-- atoms_crosslingual/      # Cross-lingual atom variants
|   |-- evaluator.py             # Atomic MQM evaluation runner
|   +-- reviews/                 # Human review annotations
|-- scripts/
|   |-- llm_generation/          # Model abstractions and generation runners
|   |   |-- models.py            # FigureAnnotator classes + MODEL_REGISTRY
|   |   |-- prompts.py           # Prompt loader (per language x figure type)
|   |   |-- prompts/             # Prompt templates (bar_chart.txt, etc.)
|   |   |-- run.py               # Main generation entry point
|   |   +-- run_adversarial.py   # Adversarial generation runner
|   |-- evaluation/              # MQM judge pipeline
|   |   |-- mqm_evaluator.py     # Core MQM scoring + judge prompts (A/B/C)
|   |   |-- judge_backends.py    # Backend router (OpenRouter/Azure/Vertex)
|   |   |-- judge_backends.json  # Backend configuration
|   |   |-- evaluate_reference_free.py
|   |   |-- evaluate_reference_only.py
|   |   |-- evaluate_reference_with_image.py
|   |   +-- export_dashboard_data.py
|   |-- experiments/             # Adversarial experiment runners
|   |   |-- run_transforms.py
|   |   |-- run_caption_bias.py
|   |   |-- run_prompt_reverse.py
|   |   |-- run_questions.py
|   |   |-- run_inductance.py
|   |   |-- run_active_admittance.py
|   |   +-- evaluate_*.py        # Corresponding evaluators
|   +-- evaluation-results/      # Result computation + plotting
|       |-- rq1/                 # Leaderboard, error analysis, heatmaps
|       |-- rq2/                 # Degradation under transforms
|       |-- rq3/                 # Capability question accuracy
|       |-- rq4/                 # Inter-judge agreement (ARI)
|       +-- discussion/          # Rank flow, scatter plots
|-- output/
|   |-- generation/              # Model outputs (per model/subfolder/fig)
|   |-- evaluation/              # Judge results (per judge_type/judge/model)
|   |-- experiments/             # Adversarial experiment outputs
|   +-- evaluation-results/      # Computed tables and plots
|-- dashboard/                   # React + Vite + Tailwind interactive UI
|   |-- src/pages/               # Page components (10 pages)
|   +-- public/                  # Exported data + images
|-- HumanEval/                   # Label Studio exports and human eval results
+-- thesis/                      # LaTeX thesis source
\end{verbatim}
}

% -----------------------------------------------------------------------------
\section{Code Architecture}
\label{sec:code-architecture}

\subsection{Model Abstraction Layer}

All VLMs are accessed through the \texttt{FigureAnnotator} abstract base class defined in \texttt{scripts/llm\_generation/models.py}. The class hierarchy is:

\begin{figure}[ht]
\centering
\begin{tikzpicture}[
  node distance=0.9cm and 1.6cm,
  class/.style={rectangle, draw, minimum width=3.2cm, minimum height=0.7cm, align=center, font=\small},
  arrow/.style={-{Stealth[length=2.5mm]}, thick}
]
  \node[class, fill=gray!15] (base) {\texttt{FigureAnnotator}\\\textit{(ABC)}};
  \node[class, fill=blue!10, below left=1.2cm and 0.5cm of base] (azure) {\texttt{AzureAnnotator}};
  \node[class, fill=green!10, below=1.2cm of base] (openrouter) {\texttt{OpenRouterAnnotator}};
  \node[class, fill=orange!10, below right=1.2cm and 0.5cm of base] (llama) {\texttt{Llama4Maverick-}\\{\texttt{Annotator}}};

  \draw[arrow] (azure) -- (base);
  \draw[arrow] (openrouter) -- (base);
  \draw[arrow] (llama) -- (base);
\end{tikzpicture}
\caption{FigureAnnotator class hierarchy.}
\label{fig:annotator-classes}
\end{figure}

Each annotator implements \texttt{annotate\_figure(prompt, image\_path, caption, paper\_title)}, which encodes the image as Base64, constructs the API request, and returns the generated description. A shared \texttt{\_retry()} wrapper handles transient failures with exponential backoff.

The \texttt{get\_annotator(model\_name)} factory function resolves a model name string to the correct annotator class using three registries:

\begin{itemize}
  \item \texttt{AZURE\_MODELS} -- maps names to Azure deployment IDs (e.g., \texttt{"gpt-5.2"} $\to$ \texttt{"gpt-5-2"})
  \item \texttt{OPENROUTER\_MODELS} -- maps names to OpenRouter model IDs
  \item \texttt{CUSTOM\_TOKEN\_MODELS} -- models requiring non-default \texttt{max\_tokens} (e.g., Gemini 3.1 Pro uses 16,000)
\end{itemize}

\subsection{Judge Pipeline}

The MQM evaluation system in \texttt{scripts/evaluation/mqm\_evaluator.py} implements three judge configurations:

\begin{enumerate}
  \item \textbf{Judge A (Reference-free)}: image + generation prompt + caption + model output. The judge evaluates accuracy against the image and completeness against the prompt.
  \item \textbf{Judge B (Reference-only)}: groundtruth annotation + model output, no image. Uses the reference as the sole source of truth.
  \item \textbf{Judge C (Reference-with-image)}: image + groundtruth + model output. The image is the primary source of truth; the reference provides expert context.
\end{enumerate}

Each judge returns a structured JSON with categorised errors, which are validated and scored using the MQM weight table:

{\scriptsize
\begin{verbatim}
  MQM_WEIGHTS = {
      ("Accuracy", "Major"): 5.0,    ("Accuracy", "Minor"): 2.0,
      ("Completeness", "Major"): 3.5, ("Completeness", "Minor"): 1.5,
      ("Clarity and Readability", "Major"): 2.0,
      ("Clarity and Readability", "Minor"): 1.0,
  }
  mqm_score = max(0, 100 - total_penalty)
\end{verbatim}
}

Judge routing is configured in \texttt{judge\_backends.json}, which maps model prefixes to backends: \texttt{azure/*} routes to Azure OpenAI, \texttt{gemini-*} to Vertex AI, and all others to OpenRouter.

% -----------------------------------------------------------------------------
\section{Data Flow}
\label{sec:data-flow}

Figure~\ref{fig:full-data-flow} shows the complete data flow from raw figures through to the dashboard.

\begin{figure}[ht]
\centering
\begin{tikzpicture}[
  node distance=1.0cm,
  box/.style={rectangle, draw, rounded corners, minimum width=3.4cm, minimum height=0.7cm, align=center, font=\footnotesize},
  data/.style={rectangle, draw, dashed, minimum width=3.0cm, minimum height=0.6cm, align=center, font=\footnotesize\itshape},
  arrow/.style={-{Stealth[length=2.5mm]}, thick}
]
  % Row 1: Inputs
  \node[data, fill=yellow!15] (figs) {Dataset/figures/\\1005 PNGs};
  \node[data, fill=yellow!15, right=2.5cm of figs] (gt) {Dataset/groundtruth/\\JSON annotations};
  \node[data, fill=yellow!15, right=2.5cm of gt] (prompts) {prompts/\\per language $\times$ type};

  % Row 2: Generation
  \node[box, fill=green!12, below=1.2cm of figs] (gen) {\texttt{run.py}\\LLM Generation};

  % Row 3: Generation output
  \node[data, fill=cyan!10, below=0.8cm of gen] (genout) {output/generation/\\model/subfolder/*.json};

  % Row 4: Evaluation
  \node[box, fill=orange!12, below=0.8cm of genout] (eval) {\texttt{evaluate\_*.py}\\MQM Judges A/B/C};

  % Row 5: Eval output
  \node[data, fill=cyan!10, below=0.8cm of eval] (evalout) {output/evaluation/\\judge\_type/judge/model/*.json};

  % Row 6: Results + Dashboard
  \node[box, fill=red!12, below left=0.8cm and -0.5cm of evalout] (compute) {\texttt{compute\_*.py}\\Result scripts};
  \node[box, fill=purple!12, below right=0.8cm and -0.5cm of evalout] (dash) {\texttt{export\_dashboard}\\React Dashboard};

  % Arrows
  \draw[arrow] (figs) -- (gen);
  \draw[arrow] (gt.south) -- ++(0,-0.35) -| (gen.north east);
  \draw[arrow] (prompts.south) -- ++(0,-0.7) -| (gen.east);
  \draw[arrow] (gen) -- (genout);
  \draw[arrow] (genout) -- (eval);
  \draw[arrow] (gt.south) -- ++(0,-0.35) -| (eval.north east);
  \draw[arrow] (figs.south) -- ++(-0.5,-0.35) |- (eval.west);
  \draw[arrow] (eval) -- (evalout);
  \draw[arrow] (evalout) -| (compute);
  \draw[arrow] (evalout) -| (dash);
\end{tikzpicture}
\caption{Complete data flow from raw figures to results and dashboard.}
\label{fig:full-data-flow}
\end{figure}

% -----------------------------------------------------------------------------
\section{Adding New Models}
\label{sec:adding-models}

To add a new VLM to the evaluation, modify \texttt{scripts/llm\_generation/models.py}:

\begin{enumerate}
  \item \textbf{For OpenRouter-hosted models}, add an entry to \texttt{OPENROUTER\_MODELS}:
{\scriptsize
\begin{verbatim}
  OPENROUTER_MODELS = {
      ...
      "new-model-name": "provider/model-id-on-openrouter",
  }
\end{verbatim}
}

  \item \textbf{For Azure-deployed models}, add to \texttt{AZURE\_MODELS}:
{\scriptsize
\begin{verbatim}
  AZURE_MODELS = {
      ...
      "new-azure-model": "deployment-name-in-azure",
  }
\end{verbatim}
}

  \item \textbf{For models with custom token limits}, add to \texttt{CUSTOM\_TOKEN\_MODELS}:
{\scriptsize
\begin{verbatim}
  CUSTOM_TOKEN_MODELS = {
      ...
      "new-model": ("provider/model-id", 16000),
  }
\end{verbatim}
}

  \item \textbf{For models requiring a new backend} (e.g., a custom API), subclass \texttt{FigureAnnotator} and implement \texttt{annotate\_figure()} and \texttt{model\_name}. Register the class in \texttt{MODEL\_REGISTRY}.
\end{enumerate}

After adding the model, re-run the pipeline:
{\scriptsize
\begin{verbatim}
  python3 scripts/llm_generation/run.py new-model-name --workers 4
  python3 scripts/evaluation/evaluate_reference_free.py new-model-name --judge azure/gpt-4.1
  python3 scripts/evaluation/evaluate_reference_only.py new-model-name --judge azure/gpt-4.1
  python3 scripts/evaluation/evaluate_reference_with_image.py new-model-name --judge azure/gpt-4.1
\end{verbatim}
}

% -----------------------------------------------------------------------------
\section{Adding New Languages}
\label{sec:adding-languages}

The system supports English, German, Bulgarian, and Chinese. To add a new language:

\begin{enumerate}
  \item \textbf{Prepare figures}: Place PNG images in \texttt{Dataset/figures/<lang>\_only/}.
  \item \textbf{Create groundtruth annotations}: Write JSON files in \texttt{Dataset/groundtruth/<lang>\_only/} following the existing schema (fields: \texttt{task\_id}, \texttt{arxiv\_id}, \texttt{paper\_title}, \texttt{paper\_language}, \texttt{figure\_type}, \texttt{caption}, \texttt{annotations}).
  \item \textbf{Write generation prompts}: Add prompt templates in \texttt{scripts/llm\_generation/prompts/cross\_language/} for the new language, covering each figure type (bar chart, line plot, pie chart, default).
  \item \textbf{Create atomic MQM decompositions}: If using the atomic evaluator, add atom JSONs in \texttt{atomic\_mqm/atoms/} for the new figures.
  \item \textbf{Update subfolder lists}: Add the new subfolder name to the \texttt{SUBFOLDERS} list in:
    \begin{itemize}
      \item \texttt{scripts/llm\_generation/run.py}
      \item \texttt{scripts/evaluation/evaluate\_reference\_free.py} (and sibling evaluators)
      \item \texttt{scripts/evaluation/export\_dashboard\_data.py}
    \end{itemize}
  \item \textbf{Update language maps}: Add entries to \texttt{SUBFOLDER\_TO\_LANGUAGE} and \texttt{LANG\_MAP} dictionaries in the evaluation and results scripts.
\end{enumerate}

% -----------------------------------------------------------------------------
\section{Configuration Reference}
\label{sec:config-reference}

Table~\ref{tab:env-vars} lists all environment variables used by the system.

\begin{table}[ht]
\centering
\caption{Environment variable reference.}
\label{tab:env-vars}
{\small
\begin{tabular}{@{}llp{5.5cm}@{}}
\toprule
\textbf{Variable} & \textbf{Required by} & \textbf{Description} \\
\midrule
\texttt{OPENROUTER\_API\_KEY} & OpenRouter models & API key for OpenRouter \\
\texttt{AZURE\_OPENAI\_ENDPOINT} & Azure models & Azure OpenAI resource URL \\
\texttt{AZURE\_OPENAI\_API\_KEY} & Azure models & Azure OpenAI API key \\
\texttt{GCP\_PROJECT} & Vertex AI / Gemini & Google Cloud project ID \\
\texttt{GOOGLE\_CLOUD\_LOCATION} & Vertex AI & GCP region (default: \texttt{us-central1}) \\
\bottomrule
\end{tabular}
}
\end{table}

Table~\ref{tab:config-files} lists key configuration files.

\begin{table}[ht]
\centering
\caption{Key configuration files.}
\label{tab:config-files}
{\small
\begin{tabular}{@{}lp{7cm}@{}}
\toprule
\textbf{File} & \textbf{Purpose} \\
\midrule
\texttt{.env} & Environment variables (API keys, endpoints) \\
\texttt{scripts/evaluation/judge\_backends.json} & Maps judge model prefixes to API backends (OpenRouter, Azure, Vertex AI) \\
\texttt{scripts/llm\_generation/models.py} & Model registry: names, router IDs, backends \\
\texttt{scripts/llm\_generation/prompts/} & Generation prompt templates per language and figure type \\
\texttt{adversarial\_experiments/samples.json} & 45-figure adversarial sample manifest \\
\texttt{atomic\_mqm/atoms/} & Atomic MQM decomposition files (120 atoms) \\
\texttt{dashboard/vite.config.ts} & Vite build configuration for the dashboard \\
\bottomrule
\end{tabular}
}
\end{table}

\noindent All API calls use \texttt{temperature=0} to ensure deterministic, reproducible outputs. The retry mechanism in both generation and evaluation uses exponential backoff with a maximum of three attempts.
